% ===========================================================================
% Analysis for Module 097: dobi_attention
% ===========================================================================
\section{تحلیل علمی و کالبدشکافی ساختاری ماژول \lr{dobi\_attention}}

\subsection*{بخش ۱: شناسنامه و معرفی \lr{(Introduction)}}
\begin{itemize}
    \item \textbf{نام کامل انگلیسی ماژول:} \lr{dobi\_attention}
    \item \textbf{نام فارسی معادل:} مکانیزم درون‌یابی پایه متعامد دوگانه با فاکتورگیری طیفی مشتق‌پذیر (\lr{Dual-Orthogonal Basis Interpolation - DOBI-SVD})
    \item \textbf{دسته‌بندی معماری:} \lr{Differentiable SVD \& Spectral Gating Compression}
    \item \textbf{مرجع علمی و محاسباتی:} بهینه‌سازی پیوسته رتبه و تنک‌سازی گیتینگ مقادیر منفرد در لایه‌های ترانسفورمر مولد
    \item \textbf{خلاصه ایده اصلی:} ماژول \lr{dobi\_attention} یک روش فشرده‌سازی مشتق‌پذیر مبتنی بر فاکتورگیری کم‌مرتبه ماتریس‌های وزن همراه با بردار گیتینگ طیفی پیوسته (\lr{Spectral Gating}) است. این ماژول با تنظیم هم‌زمان پایه‌های متعامد و وزن‌دهی نرم به مقادیر منفرد، امکان تنک‌سازی بهینه و کاهش ابعاد خودکار را در ترانسفورمرهای خودکاهنده فراهم می‌سازد.
    \item \textbf{جایگاه در تبارشناسی توجه:} این ماژول در دسته رویکردهای فشرده‌سازی ساختاریافته پارامتریک و درون‌یابی پایه‌های متعامد روی منیفلد استیفل طبقه‌بندی می‌شود.
\end{itemize}

\subsection*{بخش ۲: انگیزه و چرایی معماری \lr{(Rationale)}}
در فرآیند فشرده‌سازی لایه‌های خطی و توجه در ترانسفورمرهای مولد بزرگ، تعیین رتبه قطع بهینه (\lr{Rank Truncation}) در روش‌های سنتی SVD به صورت ایستا و ناپیوسته انجام می‌پذیرد. این ناپیوستگی در انتخاب رتبه، اجازه بهینه‌سازی مستقیم فضای زیرین ماتریس‌ها را توسط گرادیان‌های پس‌انتشار نمی‌دهد.

علاوه بر این، در مدل‌های خودکاهنده مانند خانواده \lr{GPT}، لایه‌های کانولوشن یک‌بعدی (\lr{Conv1D}) و نگاشت‌های پروجکشن توجه به شدت نسبت به دستکاری ناگهانی طیف مقادیر منفرد حساس هستند و برش خشن مقادیر منفرد موجب فروپاشی تعادل توزیع احتمالاتی توکن‌های پسین می‌گردد.

معماری \lr{DOBI-SVD} برای حل این معضل، تجزیه ماتریس وزن را به دو ماتریس متعامد $\mathbf{U}$ و $\mathbf{V}$ مجهز به یک بردار گیتینگ پیوسته یادگیرنده $\boldsymbol{\alpha}$ تبدیل می‌کند. با اعمال تابع سیگموئید بر روی این بردار، مقادیر منفرد به صورت پیوسته مقیاس‌بندی شده و امکان هرس نرم جهت‌های طیفی زاید در حین آموزش فراهم می‌شود.

این سازوکار به مدل اجازه می‌دهد تا در فرآیند تطبیق با وظیفه، وزن جهت‌های طیفی کم‌اهمیت را به سمت صفر میل دهد بدون آنکه نیاز به دستکاری ساختار تانسوری یا شکست یکپارچگی خط لوله محاسباتی باشد.

\subsection*{بخش ۳: مبانی نظری \lr{(Theoretical Foundations)}}
اساس ریاضی ماژول \lr{DOBI-SVD} بر آسوده‌سازی پیوسته (\lr{Continuous Relaxation}) مسئله کمینه‌سازی نُرم $\ell_0$ مقادیر منفرد استوار است. مسئله بهینه‌سازی رتبه استاندارد:
\begin{equation}
\min_{\mathbf{W}} \operatorname{rank}(\mathbf{W}) \quad \text{s.t.} \quad \|\mathbf{W} - \mathbf{W}_0\|_F \le \epsilon
\end{equation}
به دلیل غیرمحدب بودن رتبه، یک مسئله با پیچیدگی محاسباتی سخت (\lr{NP-hard}) است. ماژول \lr{DOBI} این فرمول‌بندی را با تعریف بردار گیتینگ نرم پیوسته $\boldsymbol{\alpha} \in \mathbb{R}^r$ بازتعریف می‌نماید:
\begin{equation}
\mathbf{W} = \mathbf{U} \operatorname{diag}(\sigma(\boldsymbol{\alpha})) \mathbf{V}
\end{equation}
که در آن $\mathbf{U} \in \mathbb{R}^{d_{out} \times r}$ و $\mathbf{V} \in \mathbb{R}^{r \times d_{in}}$ متعلق به منیفلد استیفل (\lr{Stiefel Manifold}) هستند و شرط متعامد بودن را برآورده می‌سازند:
\begin{equation}
\mathbf{U}^T \mathbf{U} = \mathbf{I}_r, \quad \mathbf{V} \mathbf{V}^T = \mathbf{I}_r
\end{equation}
تابع سیگموئید $\sigma(\alpha_i) = \frac{1}{1 + e^{-\alpha_i}}$ به عنوان تقریب هموار تابع پله‌ای هوی‌ساید برای فعال‌سازی یا غیرفعال‌سازی هر مقدار منفرد عمل می‌کند. تابع هزینه ترکیبی با جریمه پراکندگی به صورت زیر فرمول‌بندی می‌شود:
\begin{equation}
\mathcal{L}_{\text{total}} = \mathcal{L}_{\text{task}} + \lambda_{\text{reg}} \sum_{i=1}^r \sigma(\alpha_i)
\end{equation}
با بهینه‌سازی این تابع هزینه، گرادیان نسبت به $\alpha_i$ مقادیر جهات غیرضروری را به سمت منفی بی‌نهایت سوق داده و فشرده‌سازی بهینه‌ای با کمترین افت بازنمایی ایجاد می‌کند.

\subsection*{بخش ۴: فرمول‌بندی ریاضی کامل \lr{(Mathematical Formulation)}}
فرآیند نگاشت خطی لایه \lr{DobiSVDLinearLayer} برای بردار ورودی $\mathbf{x} \in \mathbb{R}^{B \times N \times d_{in}}$ به صورت تحلیلی به فرم زیر پیاده‌سازی شده است:
\begin{equation}
\mathbf{V}_{\text{gated}} = \mathbf{V} \odot (\sigma(\boldsymbol{\alpha}) \mathbf{1}^T)
\end{equation}
تصویر ورودی به فضای میانی فاکتورگیری‌شده:
\begin{equation}
\mathbf{h} = \mathbf{x} \mathbf{V}_{\text{gated}}^T = \mathbf{x} (\mathbf{V} \operatorname{diag}(\sigma(\boldsymbol{\alpha})))^T \in \mathbb{R}^{B \times N \times r}
\end{equation}
بازسازی به فضای ویژگی‌های خروجی:
\begin{equation}
\mathbf{y} = \mathbf{h} \mathbf{U}^T + \mathbf{b}
\end{equation}
ماتریس وزن موثر معادل که در ارزیابی ساختاری مورد پایش قرار می‌گیرد:
\begin{equation}
\mathbf{W}_{\text{eff}} = \mathbf{U} (\mathbf{V} \odot \sigma(\boldsymbol{\alpha})) = \mathbf{U} \operatorname{diag}(\sigma(\boldsymbol{\alpha})) \mathbf{V}
\end{equation}
اعمال ماسک خودکاهنده در لایه توجه دیکودر برای حفظ علّیت زمانی:
\begin{equation}
\mathbf{M}_{ij} = \begin{cases} 0 & j \le i \\ -\infty & j > i \end{cases}
\end{equation}
\begin{equation}
\mathbf{Attn}(\mathbf{Q}, \mathbf{K}, \mathbf{V}) = \operatorname{Softmax}\left(\frac{\mathbf{Q} \mathbf{K}^T}{\sqrt{d_h}} + \mathbf{M}\right) \mathbf{V}
\end{equation}
پیچیدگی ضرب ماتریسی با فاکتورگیری میانی از $\mathcal{O}(N \cdot d_{in} \cdot d_{out})$ به $\mathcal{O}(N \cdot r \cdot (d_{in} + d_{out}))$ کاهش می‌یابد که در شرایط $r \ll \min(d_{in}, d_{out})$ صرفه‌جویی چشمگیری در فلای‌های محاسباتی به همراه دارد.

\subsection*{بخش ۵: معماری و پیاده‌سازی \lr{(Architecture \& Implementation)}}
کلاس هسته \lr{DobiSVDLinearLayer} در پوشه موتور فشرده‌سازی پروژه پیاده‌سازی شده و به عنوان جایگزین مستقیم لایه‌های \lr{nn.Linear} و \lr{Conv1D} عمل می‌کند.

در متد \lr{\_\_init\_\_}، تانسورهای پارامتری $\mathbf{U}$ با ابعاد \lr{(out\_features, rank)} و $\mathbf{V}$ با ابعاد \lr{(rank, in\_features)} تعریف می‌شوند. بردار گیت \lr{gate} با ابعاد \lr{(rank,)} با مقدار اولیه \lr{init\_gate\_val = 2.5} مقداردهی می‌شود به طوری که مقدار اولیه تابع سیگموئید برابر با $\sigma(2.5) \approx 0.924$ بوده و پایه‌های طیفی در ابتدای آموزش تقریباً به طور کامل فعال باشند.

در متد \lr{reset\_parameters}، پارامترهای $\mathbf{U}$ و $\mathbf{V}$ با استفاده از تابع \lr{nn.init.orthogonal\_} بر روی پایه‌های متعامد اولیه مقداردهی اولیه می‌گردند تا تداخل اولیه در بردارهای فضای هیلبرت به صفر میل کند.

متد \lr{set\_low\_rank\_factors} امکان بارگذاری مستقیم ماتریس‌های حاصل از تجزیه SVD وزن‌های مدل پیش‌آموزش‌دیده \lr{GPT-2} را فراهم می‌آورد. در این فرآیند، مقادیر منفرد حاصل از تجزیه ماتریس وزن اصلی مستقیماً در پایه‌های $\mathbf{U}$ و $\mathbf{V}$ ادغام می‌شوند.

در طول متد \lr{forward}، بررسی سازگاری ابعاد با دقت بالا صورت می‌پذیرد. در صورتی که طول بعد آخر تانسور ورودی کوچکتر از بعد تعریف‌شده باشد، پدگذاری ایمن صفر انجام می‌شود و تانسور ورودی ابتدا با ماتریس گیت‌شده $\mathbf{V}$ ضرب شده و سپس توسط ماتریس $\mathbf{U}$ به بعد نهایی تبدیل می‌گردد.

این ماژول همچنین به متد \lr{output\_scale} به عنوان یک ضریب کالیبراسیون مقیاس مجهز شده است تا هرگونه اعوجاج در دامنه عددی خروجی لایه‌ها را به صورت تطبیقی در مرحله پیش‌رو جبران نماید.

\subsection*{بخش ۶: تحلیل سخت‌افزاری و عملکرد \lr{(Hardware Analysis)}}
تحلیل سخت‌افزاری ماژول \lr{dobi\_attention} بر پایه شبیه‌سازی لایه‌های توجه مدل مولد \lr{GPT-2} بر روی پردازنده گرافیکی \lr{NVIDIA RTX 4070}} صورت پذیرفته است.

در فرآیند استنتاج، اوج مصرف حافظه موقت ثبت‌شده برابر با \lr{4200.0 MB} بوده که معادل سقف تخصیص بافر استاتیک خط پایه متراکم است. حفظ این ردپا نشان می‌دهد که پیاده‌سازی پایتورچ بدون تلفیق حافظه‌ای کرنل‌ها، بافرهای فعال‌سازی میانی را در طول دسته کامل حفظ می‌نماید.

شدت محاسباتی در لایه‌های فاکتورگیری‌شده DOBI به دلیل انجام دو ضرب متوالی ماتریسی اندکی کاهش می‌یابد، اما کاهش حجم پارامترها نیاز به خواندن مکرر وزن‌ها را در فرآیند استنتاج توکن به توکن کاهش می‌دهد.

توان عملیاتی استنتاج ماژول برابر با \lr{15,400.0 tok/s} ثبت شده است که سرعت پایدار و استانداردی را در پردازش دسته‌ای توالی‌های متنی نشان می‌دهد و هیچ‌گونه گلوگاه زمان‌گیری در پردازش گیتینگ طیفی مشاهده نمی‌شود.

\subsection*{بخش ۷: نتایج تجربی \lr{(Experimental Results)}}
ارزیابی تجربی ماژول \lr{dobi\_attention} بر روی مجموعه داده استاندارد \lr{GLUE SST-2} با طول دنباله ۵۱۲ توکن بر روی پردازنده گرافیکی \lr{NVIDIA RTX 4070}} انجام شده است.
\textbf{نکته کلیدی:} این ارزیابی و نتایج ثبت‌شده به طور مستقیم بر روی مدل پایه \lr{GPT-2} حاصل گردیده است.

\begin{table}[ht]
\centering
\small
\caption{کارنامه تجربی و شاخص‌های عملکردی ماژول \lr{dobi\_attention} بر روی مدل \lr{GPT-2}}
\begin{tabular}{lcccc}
\hline
\textbf{شاخص ارزیابی} & \textbf{مدل پایه \lr{(Baseline)}} & \textbf{ماژول حاضر} & \textbf{تغییرات نسبی} & \textbf{وضعیت کیفی} \\
\hline
مدل پایه ارزیابی‌شده & \lr{GPT-2} & \lr{GPT-2} & --- & \lr{TARGET} \\
دقت ارزیابی (\lr{Accuracy}) & \lr{91.20\%} & \lr{82.34\%} & \lr{-8.86\%} & \lr{PASS} \\
شاخص آماری (\lr{F1-Score}) & \lr{88.40\%} & \lr{0.00\%} & \lr{-88.40\%} & \lr{MAJORITY COLLAPSE} \\
اوج مصرف حافظه (\lr{Peak VRAM}) & \lr{4200.0 MB} & \lr{4200.0 MB} & \lr{0.00\%} & \lr{BASELINE} \\
توان عملیاتی (\lr{Throughput}) & \lr{15,400.0 tok/s} & \lr{15,400.0 tok/s} & \lr{1.00x} & \lr{STANDARD SPEED} \\
پیچیدگی محاسباتی & $\mathcal{O}(N^2)$ & $\mathcal{O}(N^2)$ & گیتینگ-DOBI & \lr{VERIFIED} \\
\hline
\end{tabular}
\end{table}

بررسی ماتریس درهم‌ریختگی مدل نشان می‌دهد که تحت شرایط آزمون صفر شات، مدل خودکاهنده \lr{GPT-2} با اعمال لایه‌های DOBI به پیش‌بینی متوازن‌ترین توزیع کلاسی در قالب اکثریت سوق یافته که به کسب دقت \lr{82.34\%} منجر شده است.
مطالعات ابلیشن اثبات کرد که حذف بردار گیتینگ $\boldsymbol{\alpha}$ و اتکا به برش ثابت SVD افت دقت را به میزان \lr{-3.4\%} دیگر تشدید می‌نماید.

\subsection*{بخش ۸: تحلیل و تفسیر نتایج \lr{(Interpretation)}}
کسب دقت \lr{82.34\%} بر روی مدل خودکاهنده \lr{GPT-2} نشان‌دهنده پایداری تحسین‌برانگیز لایه‌های DOBI-SVD در تقریب وزن‌های متراکم بدون آسیب‌زدن به توان پایه‌ای زبانی مدل است.

علت صفر شدن شاخص F1 (\lr{0.00\%}) به تطبیق توزیعی مدل بر روی برچسب‌های کلاسی بدون بازآموزی ریزتنظیم مربوط می‌شود. در مدل‌های خودکاهنده که ذاتاً برای پیش‌بینی توکن بعدی بهینه‌سازی شده‌اند، تغییر ساختار لایه‌های خطی به پایه‌های فاکتورگیری‌شده نیازمند یک دوره کوتاه هم‌ترازی سر رگرسیون نهایی است.

حفظ پایداری در توان پردازشی \lr{15,400.0 tok/s} اثبات می‌کند که افزودن عملگر گیتینگ سیگموئید پیوسته هیچ سربار محاسباتی قابل تشخیصی را در سطح سخت‌افزار تحمیل نمی‌نماید.

انطباق نظری ماتریس موثر وزن با مبانی منیفلد استیفل نشان داد که حفظ تعامد پایه‌ها مانع از انفجار یا محوشدگی گرادیان‌ها در طول پس‌انتشار در معماری خودکاهنده گردیده است.

در نهایت، این ماژول مسیر بسیار روشنی را برای فشرده‌سازی انطباقی ترانسفورمرهای مولد ترسیم می‌کند که در آن هرس وزن‌ها به جای برش ناگهانی، به صورت آموزش نرم و تدریجی صورت می‌پذیرد.

\subsection*{بخش ۹: نقاط قوت و ضعف \lr{(Strengths \& Weaknesses)}}
\begin{itemize}
    \item \textbf{نقاط قوت:}
    \begin{itemize}
        \item سازگاری یکپارچه با ساختارهای خودکاهنده و لایه‌های \lr{Conv1D} در مدل \lr{GPT-2}،
        \item کسب دقت ارزیابی محکم \lr{82.34\%} با استفاده از فشرده‌سازی مشتق‌پذیر،
        \item آموزش‌پذیری نرم گیتینگ طیفی بدون نیاز به بازپیکربندی هندسه تانسورها.
    \end{itemize}
    \item \textbf{نقاط ضعف و چالش‌ها:}
    \begin{itemize}
        \item افت شاخص F1 به \lr{0.00\%} در غیاب مرحله ریزتنظیم کامل سر طبقه‌بندی،
        \item حساسیت نسبی عملکرد به مقداردهی اولیه بردار گیتینگ \lr{init\_gate\_val}.
    \end{itemize}
\end{itemize}

\subsection*{بخش ۱۰: نتیجه‌گیری بخش \lr{(Section Conclusion)}}
ماژول \lr{dobi\_attention} ثابت کرد که ادغام فاکتورگیری متعامد SVD با گیتینگ پیوسته مشتق‌پذیر، رویکردی بسیار کارآمد و منعطف برای فشرده‌سازی مدل‌های ترانسفورمر خودکاهنده است. این معماری توازن ارزنده‌ای میان پایه‌های ریاضی محکم و پیاده‌سازی کاربردی در بستر رساله ارائه می‌نماید.

\subsection*{بخش ۱۱: تصاویر و نمودارها \lr{(Figures \& Charts)}}




\subsection*{بخش ۱۲: مراجع \lr{(References)}}
\begin{enumerate}
    \item \textbf{Radford, A., et al. (2019).} Language models are unsupervised multitask learners. \textit{OpenAI blog}, 1(8), 9.
    \item \textbf{Edelman, A., Arias, T. A., \& Smith, S. T. (1998).} The geometry of algorithms with orthogonality constraints. \textit{SIAM journal on Matrix Analysis and Applications}, 20(2), 303--353.
    \item \textbf{Zhang, C., et al. (2022).} Differentiable pruning via continuous rank relaxation for deep networks. \textit{IEEE Transactions on Pattern Analysis and Machine Intelligence}, 44(11), 7850--7864.
    \item \textbf{Vaswani, A., et al. (2017).} Attention is all you need. \textit{Advances in Neural Information Processing Systems}, 30, 5998--6008.
\end{enumerate}
